\section{The full Hecke algebra, level, and twists}
\label{sec:hecke}

The regular point produced in Proposition~\ref{prop:regular-point}
is classical up to Tate twist, and its classical level may be divisible
by~$2$.  We now show that both features are compatible with the precise
Hecke algebra of Theorem~\ref{thm:main}.  The same argument will justify
the continuous cyclotomic twist used to normalize the original family.

For this section write $A(N)=\mathbb T_2(N)$ and
$A_k(N)=\Z_2\otimes_{\Z}T_{\leq k}^{(2)}(N)$, as in
\eqref{eq:full-hecke}. Thus an integral polynomial in the Hecke
generators tends to zero if it tends to zero on every fixed finite
range of positive weights.

We use the following structural facts at their stated scope.
The algebra $A(N)$ is reduced and is a finite product of complete
Noetherian local rings with finite residue fields; the topology just
defined induces the maximal-ideal topology on every factor.
It carries a continuous two-dimensional determinant of
$G_{\Q,S}$, where $S$ consists of the primes dividing $2N$, with
\begin{equation}\label{eq:hecke-frobenius}
 t(\operatorname{Frob}_\ell)=T_\ell,
 \qquad d(\operatorname{Frob}_\ell)=\ell S_\ell
 \quad (\ell\nmid 2N).
\end{equation}
Its classical coefficient points are Zariski dense, and the displayed
Hecke elements topologically generate it.  These are
\cite[Lemmas~2.1--2.2 and Proposition~2.3]{Dim}; the Frobenius
convention is exactly that of Equation~(2.2) there.  In particular, the
characteristic polynomial in this convention is
$Z^2-T_\ell Z+\ell S_\ell$.  Finally,
\cite[Theorem~1.1 and Proposition~3.1]{Dim} say, respectively, that
every irreducible component has dimension four, that the closure of
bounded-weight classical points has dimension at most one, and that
the Eisenstein closure has dimension at most two.  It follows that
classical cuspidal points of weight at least three are dense on every
component: remove the latter two closed loci and the other components,
and apply classical density in every nonempty remaining open subset.

These facts remain valid after finite extension of integral
coefficients and passage to a residue factor.  Indeed this extension
is finite flat.  After inverting $2$ it is separable, so reducedness
persists; integral torsion-freeness then gives reducedness before
inverting $2$.  Testing in all coefficient embeddings preserves the
stated density.  Integral finite extensions preserve component
dimensions.  We will use these observations when a character or an
eigenpacket requires larger coefficients.

\subsection{Evaluation on the ordinary tower}

\begin{proposition}\label{prop:hecke-level}
Let $N\geq5$ be odd, and let $f$ be a classical cuspidal eigenform
of positive weight and level dividing $2^aN$, for some $a\geq0$.
Its eigenvalues away from $2N$ define a continuous homomorphism
\[
 A(N)\longrightarrow\OO_{E_f}
\]
after choosing a finite $2$-adic coefficient field $E_f$ containing
them.  The associated determinant has the Frobenius polynomials
of~$f$ in \eqref{eq:hecke-frobenius}.
\end{proposition}

\begin{proof}
We evaluate all the forms on one ordinary tower.  Density of the
tame-level forms on that tower will then turn boundedness of the
Hecke operators into the continuity required here.

\emph{The ordinary tower.}  Enlarge an
integer coefficient ring $\OO$ as necessary.  Over the formal ordinary
locus of the compactified tame modular curve of level $\Gamma_1(N)$,
consider the tower of trivializations
\[
 \iota:\widehat{\mathbb G}_m\xrightarrow{\sim}\widehat E
\]
of the formal group of the universal generalized elliptic curve.
The ordinary locus here includes the cusps.  Write $\mathcal I$ for
the ring of integral functions on this tower, completed for the
coefficient-uniformizer topology, and put
$\mathcal B=\mathcal I[1/2]$, with unit ball $\mathcal I$.
The tower is the inverse system of finite \'etale trivialization
torsors of the ordinary connected $2$-divisible group.  One can see
this by Cartier duality: that group's dual is \'etale of height one.
This also proves the assertion at a multiplicative cusp.
In particular, the tower and $\mathcal I$ are flat over $\OO$.
We use the ordinary-tower construction of
\cite[Section~4.2, Lemma~4.2.1 and Theorem~4.2.2]{Katz1973}.
For Drinfeld level structures we use
\cite[Section~3.2 and Lemma~3.5.1]{KatzMazur1985}; the compactified
moduli interpretation, including cusps in bad characteristic, is
\cite[Theorem~1.2.1 and Definition~2.4.1]{Conrad2007}.

A section of the $i$th power of the Hodge line $\omega$ evaluates
to a function on this tower by expressing its pullback using the
standard differential $dT/(1+T)$ on $\widehat{\mathbb G}_m$.

\emph{Density.} We claim that the $\operatorname{Frac}(\OO)$-linear
span of the evaluations of tame-level classical forms of positive
weights is dense in $\mathcal B$.  Only rational density is needed.

Here are the details at the prime two.  Let $H$ be the weight-one
Hasse invariant on the residue curve.  Its evaluation on the tower
is $1$, and the normalized weight-four Eisenstein series $E_4$
reduces to $H^4$; see \cite[Section~2.1]{Katz1973}.
The open set where $E_4$ is nonvanishing is affine, since $\omega$
is ample.  Its formal completion is the ordinary base; denote its
ring of functions by $\mathcal I_0$.
Serre vanishing for a sufficiently high power of $\omega$ lifts
$H^{4m+1}$ to an integral form $G$ of weight $4m+1$.
Consequently
\[
 \eta=G/E_4^m
\]
is a basis of $\omega$ on the formal ordinary base.  Choose a
parameter $z$ on the formal group whose cotangent is $\eta$.
Such a parameter is obtained by lifting this cotangent successively
on the affine base.  Express the universal trivialization as
\[
 \iota^*z=a_1T+a_2T^2+\cdots.
\]
The ring $\mathcal I$ is topologically generated over $\mathcal I_0$
by the coefficients $a_j$ and $a_1^{-1}$: formal group isomorphisms
are represented by their coefficients, the formal group identities,
and invertibility of the first coefficient. Over the ordinary base
modulo $\varpi^h$, the ideals cutting out the $2^n$-torsion are cofinal
with the powers of the parameter ideal: the formal group has height
one and $2$ is nilpotent on this base. Thus an isomorphism of the formal
groups is a compatible system of isomorphisms of their finite flat
$2^n$-torsion groups. Cartier duality identifies this functor with the
inverse limit of the finite \'etale trivialization torsors. The
coefficient presentation and the tower therefore have the same
coordinate ring modulo every $\varpi^h$. After inverting $2$, formal logarithms
express each individual $a_j$ as a polynomial in $a_1$ with
coefficients in $\mathcal I_0[1/2]$.

Let $\mathcal C$ be the closed linear span under consideration.
It is an algebra, since multiplication adds positive weights.
The evaluations of $E_4$ and $G$ are $1$ modulo the uniformizer.
Their powers with exponents $2^n$ tend to $1$, so $1\in\mathcal C$;
their powers with exponents $2^n-1$ then show that their inverses
also belong to $\mathcal C$.
Ratios of equal-weight sections by powers of $E_4$ give all
functions on the affine base, by the section-ring description of
an ample line bundle.  Completion gives
$\mathcal I_0[1/2]\subseteq\mathcal C$.
Since $\eta=G/E_4^m$ evaluates to $a_1$, both $a_1$ and
its inverse belong to $\mathcal C$.  The logarithm identities now give every
$a_j\in\mathcal C$, and topological generation proves
$\mathcal C=\mathcal B$.
No uniform bound on the denominators of all logarithm coefficients
is required: each coefficient is handled separately before taking
the closure.

\emph{Dyadic level and Hecke operators.} We next evaluate the given form on the same tower and check
that its Hecke operators preserve the integral unit ball.
For forms of level $\Gamma_1(2^aN)$, choose a primitive $2^a$th root of unity in the coefficient ring
and transport its multiplicative torsion point by $\iota$.
This supplies the extra Drinfeld $\Gamma_1(2^a)$-structure.
The construction works on generalized elliptic curves at cusps:
the original tame level is ample, and adding a subgroup in the
identity component preserves ampleness.  The formal group of a
Tate polygon is $\widehat{\mathbb G}_m$ by
\cite[Equation~(2.5.3)]{Conrad2007}.  The integral sections of the
Hodge line recover classical forms after inverting $2$
\cite[Section~4.4, Equation~(4.4.2)]{Conrad2007}; hence a fixed
scalar multiple of a classical form extends over the integral
model and has evaluation in $\mathcal I$.  Evaluation is nonzero for a
nonzero form, as its $q$-expansion at a compatible multiplicative
cusp shows.

The good Hecke operators preserve $\mathcal I$ and commute with
all these evaluations.  For $T_\ell$ the correspondence sums over
degree-$\ell$ isogenies, with normalization $\ell^{-1}$, a unit
in $\OO$.  The isogenies induce isomorphisms of formal groups,
and the finite flat correspondence and its trace are integral.
The extension over multiplicative cusps, and the normalization
of trace--pullback as $\ell T_\ell$, are
\cite[Theorem~4.4.3 and Equation~(4.5.1)]{Conrad2007}.
For $\ell S_\ell$, the diamond operator and the scalar change
of trivialization contribute $\ell^i\ell^{-1}$ in weight $i$,
which is the required $\ell^{i-1}\langle\ell\rangle$.
This includes the diamond action on the additional $2$-power
level.  Thus every integral polynomial in these operators is
a contraction of $\mathcal B$.

\emph{Continuity.} Let $h_\alpha$ be a net of such polynomials
that tends to zero at each finite weight stage.  On every finite
sum of evaluated tame classical forms, $h_\alpha$ tends to zero.
Density and the uniform contraction bound imply
\[
 h_\alpha v\longrightarrow0\qquad(v\in\mathcal B):
\]
first approximate $v$ by a finite sum, then use the contraction
bound on the error.  Apply this to the nonzero evaluation $v_f$
of~$f$.  Since $h_\alpha v_f=\lambda_f(h_\alpha)v_f$,
the eigenvalues tend to zero.  The same argument shows that a
polynomial relation in $A(N)$ acts as zero on $v_f$.
The eigenvalue map therefore extends continuously from the dense
polynomial algebra to $A(N)$.  Its values are integral, either
by the contraction bound or by classical integrality of good
Hecke eigenvalues.  Compatibility of evaluation with the Hecke
action proves the asserted Frobenius identities.
\end{proof}

\subsection{Continuous cyclotomic twists}

\begin{proposition}\label{prop:hecke-twists}
Let $N\geq5$ be odd.  Let $\lambda:A(N)\to\OO_E$ be a continuous
coefficient point, with associated two-dimensional determinant
$(t_\lambda,d_\lambda)$.  If $E'/E$ is finite and
$\theta:\Z_2^\times\to\OO_{E'}^\times$ is continuous, then
the twisted determinant
\[
 \bigl((\theta\circ\eps)t_\lambda,
       (\theta\circ\eps)^2d_\lambda\bigr)
\]
is also supplied by a continuous coefficient point of $A(N)$.
In particular this holds for every integer Tate twist, with the
same tame level $N$.
\end{proposition}

\begin{proof}
Fix integral coefficients containing the residual values, and
work in the factor $A$ of the resulting Hecke algebra selected
by $\lambda$.  Let $R^{\mathrm{ps}}$ be the universal global
two-dimensional determinant deformation ring for these residual
data, with ramification allowed at the primes dividing $2N$ and
with varying determinant.  The structural facts above give a
continuous surjection
\[
 R^{\mathrm{ps}}\twoheadrightarrow A,
 \qquad J=\ker(R^{\mathrm{ps}}\longrightarrow A).
\]

First suppose $\theta$ has finite order, and enlarge coefficients
to contain its values.  Twisting the universal determinant over
$A$ defines a continuous map $R^{\mathrm{ps}}\to A$ over these
enlarged coefficients.  There is no change of residual factor:
every finite quotient of $\Z_2^\times$ is a $2$-group, so a
finite-order character reduces to $1$ in characteristic two.
At every classical cuspidal point of weight at least three,
the twisted determinant belongs to a classical finite-character
twist.  Its extra conductor is a power of $2$; Proposition~\ref{prop:hecke-level}
therefore supplies a coefficient point of the same $A(N)$.
For $j\in J$, its image under the twisted map vanishes at all
these points.  Their density and reducedness show that this image
is zero.  Thus twisting factors through $A$ and preserves all its
coefficient points.  Equality of the specialized determinants
follows from \eqref{eq:hecke-frobenius} and Chebotarev.

Now fix $\lambda$, and fix a sign $s\in\{1,-1\}$.  Every
$u\in\Z_2^\times$ has a unique expression $u=(-1)^e5^b$,
with $e\in\{0,1\}$ and $b\in\Z_2$.  Define the universal character
\[
 \Theta_X((-1)^e5^b)=s^e(1+X)^b
 \quad\text{in }\OO_{E'}[[X]]^\times.
\]
Twisting the determinant of $\lambda$ by $\Theta_X\circ\eps$
gives a continuous map
$R^{\mathrm{ps}}\to\OO_{E'}[[X]]$.
For every $2$-power root of unity $\zeta$, evaluation at
$X=\zeta-1$ is a finite-order twist, so the image of every
$j\in J$ vanishes at all such values.  A nonzero series in
$\OO_{E'}[[X]]$ cannot do this.  Indeed, after dividing by the
largest common uniformizer power of its coefficients,
Weierstrass preparation expresses it as a distinguished polynomial
times a unit.  The polynomial has only finitely many roots in
the open unit disk.  Therefore the image of $J$ is zero, and the
universal twisted map factors continuously through $A$.

For the given continuous $\theta$, its image on the pro-$2$
group $1+4\Z_2$ reduces trivially in the odd-order group
$\kk_{E'}^\times$.  Hence $\theta(5)-1$ lies in the maximal
ideal of $\OO_{E'}$.  Specialize $s=\theta(-1)$ and
$X=\theta(5)-1$.  This continuous specialization gives the
required point.  All ring topologies used here are the completed
local topologies, which agree with the prescribed finite-stage
topology by \cite[Proposition~2.3]{Dim}.
\end{proof}

We can now pass from classical representations in the companion
normalization to points of $A(N)$.  If $V_f$ denotes the arithmetic representation of a
weight-$k$ primitive form with nebentype $\nu$, then
$\det V_f=\nu\eps^{k-1}$, and the wedge pairing gives
\begin{equation}\label{eq:classical-dual-conversion}
 V_f^\vee\simeq V_f\otimes\nu^{-1}\eps^{1-k}
       \simeq V_g\otimes\eps^{1-k},
\end{equation}
where $g$ is the primitive form associated with
$f\otimes\nu^{-1}$.  Thus an expression
$r_x\simeq V_f^\vee\eps^n$ becomes
$r_x\simeq V_g\eps^{n+1-k}$.
This is the conversion in \cite[Section~2.1]{FM2}.
Propositions~\ref{prop:hecke-level} and~\ref{prop:hecke-twists}
therefore put every representation classical up to Tate twist
on $A(N)$ for a sufficiently divisible odd $N$.
